\subsection{Product Sampling in Constant Rounds}\label{sec:rounds}

\paragraph{Achieving constant rounds through parallel repetition.}
In Sections~\ref{sec:prod}, we showed a sublinear communication protocol
for product sampling when $\ipw$ is sufficiently large.  Moreover, this
protocol provably leaked no more information than the inner product.
However, this protocol required $O(1/\ipw)$ rounds of communication.
This raises the question of whether constant-round sublinear product
sampling is possible under the same restrictions on the inputs.

Our protocol to achieve this takes a relatively standard approach.
Suppose that we are given the value of $\ipw$.  Then, since the expected
number of samples until a collision is a function of $\ipw$, we can just
run the inner loop of protocol $\Pp$ in parallel sufficiently many times
to guarantee that the protocol would terminate with all but negligible
probability.  

\paragraph{How many times to repeat?} However, there is one catch.  It
is not actually possible to compute $\ipw$ in sublinear communication!
One simple solution is to use our promise on the input:~we could run the
inner loop enough times to guarantee termination for any inputs
satisfying the promise (e.g.~$\omega(\frac{n}{\log n}$) times).
However, this forces us to adopt the worst-case communication cost,
which might be undesirable.  (Recall, it also offers the least leakage,
which might be desirable.) 
%
Instead, we re-establish the trade-off between leakage and efficiency as
follows.  We begin by computing an approximation of the inner product in
sublinear communication (see Section~\ref{sec:jlt}). Using this
approximation, we can then realize our sublinear communication, constant
round protocol for product sampling as follows in the next subsection.  

\subsubsection{Secure approximation of the inner product}
\label{sec:jlt}
We achieve a protocol that securely approximates the inner product with
sublinear communication. In particular, we take advantage of the well
known Johnson–Lindenstrauss Transform (JLT)~\cite{JL84,STOC:IndMot98}
sketch. 


\paragraph{Additional assumptions about $\vec{w_1}$ and $\vec{w_2}$}. 
We assumed that $\vec{w_1}$ and $\vec{w_2}$ are normalized and correlated such
that $\langle \vec{{w}_1},{\vec{w}_2} \rangle = \omega(\log n/n)$. In a similar vein, we assume that the
cosine similarity of the two vectors $\w_1$ and $\w_2$ is not small,
e.g., $\omega(1/\log n)$. 

Recall the cosine similarity between the two vectors $\vec{w}_1$ and
$\vec{w}_2$ is defined as %
$\cos(\w_1, \w_2) = \frac{\langle \vec{w}_1, \vec{w}_2 \rangle}{ \|
\vec{w}_1 \|_2 \cdot \| \vec{w}_2 \|_2}$.
% 
Since the $L_1$ norm of each vector is equal to 1, their $L_2$ norms
will typically much smaller than 1, which implies that the cosine
similarity is usually much larger than $\ipw$. 


\paragraph{Approximating the inner product using JLT sketches.}
The JLT sketch of $\vec{x}$ is equal to $\vec{M} \vec{x}$, where $\vec{M}$ is a
random $k \times n$ matrix with $k \ll n$. More specifically, the inner product
of the two vectors is approximated as follows:

\BEN
\item[] \hspace{-2em} ${\sf approxIP}(\w_1, \w_2)$: 
	~~~ $\rhd$ $\w_1$ and $\w_1$ are $n$ dimensional vectors. 
\item Choose $k \times n$ matrix $\vec{M}$ such that each entry $M_{i,j}$ is chosen from an
independent Gaussian distribution of mean $0$ and variance $1$.

\item Output $\frac 1 k \cdot \langle \vec{M}\vec{w}_1, \vec{M}\vec{w}_2 \rangle$. (Here, we
slightly abuse the notation and treat the vectors $\w_1$ and $\w_2$ as column vectors.)
\EEN

\begin{lemma} {\rm (cf.~\cite[Corollary 3.1]{Kaban15})}
\label{lem:jlt_k}
For all $\w_1, \w_2$ such that $\cos(\w_1, \w_2)\geq t$, the
procedure ${\sf approxIP}(\w_1, \w_2)$ approximates $\langle \vec{w}_1,
\vec{w}_2 \rangle$ up to a $1\pm \epsilon$ approximation factor with all but
negligible probability (over the choice of the JLT matrix), using JLT dimension
$k = \omega\left(\frac{\log(n)}{t^2 \cdot \epsilon^2}\right)$. 
\end{lemma}


\paragraph{Privacy of the approximate output.} What is interesting is that the
approximate inner product doesn't reveal anything more than the inner product
itself. In this sense, it satisfies the notion of private
approximation introduced in~\cite{FIM+06}. In particular, we prove the following:

\begin{lemma}\label{lem:simip}
The output of $\mathsf{approxIP}(\vec{w}_1, \vec{w}_2)$ can be simulated
perfectly given only $\langle \vec{w}_1, \vec{w}_1 \rangle$, $\langle
\vec{w}_2, \vec{w}_2 \rangle$, and $\langle \vec{w}_1, \vec{w}_2 \rangle$.
\end{lemma}

The proof is found in Appendix~\ref{apx:pfLemma10}. 


\paragraph{Private protocol via JLT.} 
Using the JLT sketch, we can design a private protocol approximating the
inner product. See Protocol~\ref{fig:ip-fhe}. The protocol uses
threshold FHE (e.g., ~\cite{popets:MTBH21}). 
%and (for simplicity) works
%in the pre-processing model.  We note, however, that this pre-processing
%could be elminated by having the parties jointly generate an encrypted
%random seed $s$ and using randomness generated by applying a PRG to the
%seed $s$ (under FHE encryption) to sample the JLT matrix $\lbra \vec M
%\lket$.  \dnote{Should the PRG solution be emphasized more? B/c some
%people may not like pre-processing. We already require FHE (and
%essentially arbitrary computation under it) so not sure what the
%drawback is of the PRG-based solution.} \anote{I agree with Dana here.
%I think we should present the PRG-based solution instead of the
%pre-processing model.}

\floatname{algorithm}{Protocol}
\begin{algorithm}[htbp]



\paragraph{Inputs:} Parties $P_1$ and $P_2$ has inputs $\w_1$ and $\w_2$
respectively. 

\medskip

The protocol proceeds as follows:
\BEN
    \item Parties set up a threshold FHE scheme.  \item They securely
    sample $k \times n$ matrix $\vec M$ described in the above with in
    the threshold FHE. In particular, they jointly generate an encrypted
    random seed $\lbra s \lket$. Using this randomness, parties
    homomorphically evaluates $\lbra PRG(s) \lket$, where $PRG$ is a
    pseudorandom generator, to obtain the JLT matrix $\lbra \vec M
    \lket$.
    
    \item Each party  $P_b$ homomorphically evaluates $\lbra \tilde \w_b \lket =
    \lbra \vec M \w_b \lket$.

    \item Party $P_1$ sends $\lbra \tilde \w_1 \lket$ to $P_2$.  
    
    \item Party $P_2$ homomorphically evaluates $\lbra \langle \tilde
    \w_1, \tilde \w_2 \rangle \lket$ and sends it to  $P_1$. 

    \item Parties execute threshold decryption to obtain and output
    $\frac 1 k \cdot \langle \tilde
    \w_1, \tilde \w_2 \rangle$.
\EEN
\caption{Private protocol for computing approximate inner product}
\label{fig:ip-fhe}
\end{algorithm}


\paragraph{Security.}
Since every protocol message is a ciphertext, based on semantic security
of the threshold FHE, it is easy to see that the protocol securely
realizes a functionality for computing ${\sf approxIP}$. 
%
Based on Lemma~\ref{lem:simip}, the leakage profile of the functionality is
$\langle \vec{w}_1, \vec{w}_1 \rangle$, $\langle
\vec{w}_2, \vec{w}_2 \rangle$, and $\langle \vec{w}_1, \vec{w}_2 \rangle$.




\subsubsection{Constant-round protocol for product sampling}
\label{sec:const}
Note that the Protocol~\ref{prot:prodsampapx} has the following 
structure.  In particular:

\BI
\item The probability that Protocol~\ref{prot:prodsampapx} samples a
good index and halts in a given trial is $p = \bra \w_1, \w_2 \ket$.
\EI

We need to repeat $r$ trials in parallel so that the probability that
all $r$ trials fail is negligible. In other words, we should have 
%
$$(1 - p)^{r} \le e^{-p \cdot r} \le e^{-\omega(\log \secparam)}.$$ 
%
This means that we should have
$r > \frac {\omega(\log \secparam)}{p}$. 

Moreover, in the previous subsection, we discussed how to obtain a good estimate
$\tilde p = (1 \pm \e) p$. 
Therefore, we should have 
$$r > \frac {(1+\eps) \cdot \omega(\log \secparam) }{\tilde p}> \frac
{\omega(\log \secparam)} p .$$

In summary, by running $\frac {(1+\eps) \cdot \omega(\log \secparam)
}{\tilde p}$ instances in parallel, we achieve constant round protocols
for product sampling  with negligible failure probability. The final
protocol should perform extra steps to hide from which trial the output
comes from, and these changes can made in a straightforward way.



%
%We now describe a constant round protocol for product sampling given an
%$(\epsilon,\delta)$-approximation of $\ipw$.  This protocol is in the
%$\F_\oslone$-hybrid model.
%
%\floatname{algorithm}{Protocol}\label{prot:prodsampdp}
%\begin{algorithm}[htbp]
%\textbf{Inputs:} Party $P_i$ has input $\vec{w}_i$ of length $n$.   
%
%\begin{enumerate}
%    \item Parties hold $aIP$ an $(\epsilon, \delta)$-approximation of $\ip{\vec{w}_1}{\vec{w}_2}$ with $\delta=\negl(\secparam)$.
%    
%    \item The parties execute the following steps $m=\left(\frac{aIP}{1-\epsilon}\right)\cdot \omega(\log{\secparam})$ times in parallel.
%    
%    \BEN
%    \item Execute $\F_\oslone$ with $P_1$ as the sender with input $\vec{w}_1$ and $P_2$ as the receiver.  Let $i_{1,1}^j$ and $i_{1,2}^j$ be the output of the $j$th execution to $P_1$ and $P_2$ respectively. 
%    
%    \item Execute $\F_\oslone$ with $P_2$ as the sender with input $\vec{w}_2$ and $P_1$ as the receiver.  Let $i_{2,1}^j$ and $i_{2,2}^j$ be the output of the $j$th execution to $P_1$ and $P_2$ respectively.
%    
%    \EEN
%    
%    \item Execute $\F_{2PC}$ for the following circuit: 
%      \BEN
%      \item[] \hspace{-2em} Input: $(i_{1,1}^j, i_{2,1}^j, i_{1,2}^j, i_{2,2}^j)$ for $j = 1, \ldots, m$. 
%      \item Let $i_1^j = i_{1,1}^j \xor i_{1,2}^j$, $i_2^j = i_{2,1}^j \xor i_{2,2}^j$. 
%      \item Find the smallest $j$ such that $i_1^j == i_2^j$, and output $i_1^j$ to both $P_1$ and $P_2$.  If no such $j$ exists, output $\mathsf{abort}$.
%      \EEN
%
%\end{enumerate}
%\textbf{Output:} Both parties output the sampled index $i$ or $\mathsf{abort}$. 
%
%\caption{Constant-round product sampling ($\Pp^{cr}$)}
%\end{algorithm}
%
%\paragraph{Security}
%\anote{It actually seems that we no longer need to hide the round of abort?  Could you just use a comparison protocol that outputs result in the clear?}
%We will prove the following theorem.
%\begin{theorem}\label{thm:prodDP}
%$\Pp^{DP}$ is secure \anote{Need to capture leakage from approximation algorithm.}
%\end{theorem}
%Before proving the main theorem, we prove the following claim to argue that this protocol does not output $\mathsf{abort}$, except with negligible probability.
%
%\begin{claim}\label{lem:prodAbort}
%For any input vectors $\vec{w}_1$, $\vec{w}_2$ such that $\ip{\vec{w}_1}{\vec{w}_2} > 0$, the probability that Protocol~\ref{prot:prodsampdp} aborts is negligible in $\secparam$.
%\end{claim}
%\begin{proof}
%We first consider the probability that there is no colliding indices in the first $r-1$ rounds
%\begin{align*}
%\Pr[\mbox{no collision in first }r-1 \mbox{ rounds}] = (1- \ip{\vec{w}_1}{\vec{w}_2})^{r-1}
%\end{align*}
%
%\noindent
%Since $aIP$ is an $(\eps,\delta)$-approximation of $\ip{\vec{w}_1}{\vec{w}_2}$, we know that with all but negligible probability $(1-\eps)\ip{\vec{w}_1}{\vec{w}_2} \le aIP \le (1+\eps)\ip{\vec{w}_1}{\vec{w}_2}$.
%
%\medskip\noindent
%Letting $x=1/\ip{\vec{w}_1}{\vec{w}_2}$ and using the inequality $(1-1/x)^x \le e^{-1}$, we see that after $\omega(x\log{\secparam})$ rounds $\Pr[\mathsf{abort}] \le e^{-\omega(\log{\secparam})} \le \negl(\secparam)$.  Thus, given $IP$, this guarantee holds after at most $\frac{IP}{1-\eps}\cdot\omega(\log{\secparam})$ iterations.
%\end{proof}
%
%We can now prove Theorem~\ref{thm:prodDP}. Due to Claim~\ref{lem:prodAbort}, for the remainder of this proof, we assume that the protocol does not output $\mathsf{abort}$.  
%
%\begin{proof}
%The rest of the proof is essentially identical to the proof of Theorem~\ref{thm:prodIP}.  
%\end{proof}
%
%\paragraph{Performance}
%This protocol executes at most $m=\left(\frac{aIP}{1-\epsilon}\right)\cdot \omega(\log{\secparam})$ executions of the 2-PC comparison protocol in parallel, each requiring $O(\ell)$ communication.  Finally, the 2-PC for selecting the output can be done in communication linear in the number of executions, thus giving a total communication of $O(ml) = o(n)$.  As long as we use constant-round 2-PC protocols to realize all these computations, the resulting protocol also has constant rounds.
%
%
